\documentclass[11pt,a4paper]{article}
\usepackage[margin=1.8cm]{geometry}
\usepackage{amsmath,amssymb,array,tabularx,fontspec,cancel,tikz,graphicx}
\setmainfont{FreeSerif}
\setlength{\parindent}{0pt}
\setlength{\parskip}{0pt}
\newcounter{sq}
% \sq{text}: auto-numbered (Roman numerals like the original paper)
\newcommand{\sq}[1]{\stepcounter{sq}\par\medskip\noindent\hangindent=8.5em\hangafter=1
  \makebox[8.5em][l]{\textbf{\Roman{sq}.}}#1\par}

% \fd{left labels}{right labels}{edges}: arrow (mapping) diagram
\newcommand{\fd}[3]{\par\smallskip\noindent\hspace*{8.5em}\begin{tikzpicture}[>=stealth,scale=0.75]
 \draw (0,0) ellipse (0.9 and 1.5); \draw (3.5,0) ellipse (0.9 and 1.5);
 \foreach \l [count=\i] in {#1} {\node[font=\small] (L\i) at (0,1.25-0.7*\i+0.7) {\l};}
 \foreach \r [count=\i] in {#2} {\node[font=\small] (R\i) at (3.5,1.25-0.7*\i+0.7) {\r};}
 #3
 \end{tikzpicture}\par}
% \tfig{tikz code}: geometry figure
\newcommand{\tfig}[1]{\par\smallskip\noindent\hspace*{8.5em}\begin{tikzpicture}[scale=0.8,font=\small]#1\end{tikzpicture}\par}
% \pfig{file}: figure copied (cropped) from the original paper
\newcommand{\pfig}[2]{\par\smallskip\noindent\hspace*{8.5em}\includegraphics[width=#2]{#1}\par}
% score-distribution table used in several questions
\newcommand{\dist}{\par\smallskip\noindent\hspace*{8.5em}\begin{tabular}{|l|c|c|c|c|c|c|c|}\hline
\textbf{Scores} & 24-28 & 29-33 & 34-38 & 39-43 & 44-48 & 49-53 & \textbf{Total}\\ \hline
\textbf{No. of students} & 3 & 6 & 12 & 23 & 15 & 6 & 65\\ \hline\end{tabular}\par\smallskip\hspace*{8.5em}}
% biology marks tables (6 columns / 5 columns as in the original)
\newcommand{\bioA}{\par\smallskip\noindent\hspace*{8.5em}\begin{tabular}{|l|c|c|c|c|c|c|}\hline
\textbf{Marks} & 20-24 & 25-29 & 30-34 & 35-39 & 40-44 & 45-49\\ \hline
\textbf{No. of students} & 05 & 08 & 12 & 15 & 03 & 02\\ \hline\end{tabular}\par\smallskip\hspace*{8.5em}}
\newcommand{\bioB}{\par\smallskip\noindent\hspace*{8.5em}\begin{tabular}{|l|c|c|c|c|c|}\hline
\textbf{Marks} & 20-24 & 25-29 & 30-34 & 35-39 & 40-44\\ \hline
\textbf{No. of students} & 05 & 08 & 12 & 15 & 03\\ \hline\end{tabular}\par\smallskip\hspace*{8.5em}}

\begin{document}
\begin{center}
  {\Huge\bfseries The Hope Science Academy}\\[3pt]
  {\small Excellence is our standard not goal}\\[2pt]
  {\small 247-E-1, (Opposite LDA School) Johar Town, Lahore. PH: 0300-8194789, 0322-4157001}
\end{center}
\medskip\noindent
\begin{tabularx}{\linewidth}{|X|X|X|X|}
\hline
\textbf{Student Name:} & \textbf{Roll Num:} & \textbf{Class:} 10th & \textbf{Paper Code:} \\ \hline
\textbf{Subject Name:} Mathematics. & \textbf{Time Allowed:} & \textbf{Total Marks:} 824 & \textbf{Exam Date:} 09-Oct-2026 \\ \hline
\multicolumn{4}{|l|}{\textbf{Exam Syllabus:} FULL BOOK} \\ \hline
\end{tabularx}
\bigskip
\noindent\textbf{Q3. Write short answers of the following questions.}\hfill\textbf{1$\times$604=604}

\sq{Simplify the following: $i^{5}$}
\sq{Simplify the following: $i^{16}$}
\sq{Simplify the following: $(-i)^{-19}$}
\sq{Simplify the following: $27i^{-26}$}
\sq{Simplify the following: $i^{11}+i^{5}$}
\sq{Simplify the following: $(i^{4}+i^{3}+i^{2}+i)^{2}$}
\sq{Simplify the following: $\left(\dfrac{i^{8}}{i^{5}}\right)^{-5}$}
\sq{Simplify the following: $i^{13}\times i^{29}$}
\sq{Write in term of $i$: $2+\sqrt{-4}$}
\sq{Write in term of $i$: $3-\sqrt{-7}$}
\sq{Write in term of $i$: $\dfrac{2}{5}+\dfrac{\sqrt{-16}}{5}$}
\sq{Write in term of $i$: $\sqrt{2}-\sqrt{3}$}
\sq{Find the values of x and y: $(2x+5)+(y-3)i=1+2i$}
\sq{Simplify the following: $i^{7}$}
\sq{Simplify the following: $i^{8}$}
\sq{Simplify the following: $i^{17}$}
\sq{Simplify the following: $i^{-25}$}
\sq{If $(2x-1)+(y+4)i=5+7i$, then find the values of x and y.}
\sq{Simplify and write in the form of a+b$i$: $(2+5i)+(3-zi)$}
\sq{Simplify and write in the form of a+b$i$: $(16-3i)+(9+2i)$}
\sq{Simplify and write in the form of a+b$i$: $(9-2i)-(7-3i)$}
\sq{Simplify and write in the form of a+b$i$: $(11+9i)-(9-7i)$}
\sq{Simplify and write in the form of a+b$i$: $(3+4i)(2-3i)$}
\sq{Simplify and write in the form of a+b$i$: $(5-2i)(3-4i)$}
\sq{Simplify and write in the form of a+b$i$: $(3-5i)\div(2-4i)$}
\sq{Simplify and write in the form of a+b$i$: $(5+2i)\div(6-3i)$}
\sq{Write additive inverse for the following complex number: $3+2i$}
\sq{Write additive inverse for the following complex number: $4-3i$}
\sq{Write additive inverse for the following complex number: $5-7i$}
\sq{Write additive inverse for the following complex number: $-\dfrac{2}{3}+\dfrac{5}{4}i$}
\sq{Find the multiplicative inverse for following complex number: $4+5i$}
\sq{Find the multiplicative inverse for following complex number: $6+2i$}
\sq{Find the multiplicative inverse for following complex number: $7-3i$}
\sq{Find the multiplicative inverse for following complex number: $\sqrt{5}-4i$}
\sq{If $z_{1}=2+5i,\ z_{2}=1-3i$ and $z_{3}=2+i$, then verify that:\\ $z_{1}+z_{2}=z_{2}+z_{1}$}
\sq{If $z_{1}=2+5i,\ z_{2}=1-3i$ and $z_{3}=2+i$, then verify that:\\ $z_{1}z_{2}=z_{2}z_{1}$}
\sq{If $z_{1}=2+5i,\ z_{2}=1-3i$ and $z_{3}=2+i$, then verify that:\\ $z_{1}+(-z_{1})=(-z_{1})+z_{1}=0$}
\sq{Find the values of a and b, if $(a+bi)(1+3i)=-8+11i$.}
\sq{Add $(3+4i)$ and $(5-2i)$.}
\sq{Simplify: $(8+2i)-(5-6i)$}
\sq{Simplify: $(4-3i)-(2-5i)$}
\sq{Simplify: $(3-4i)(5-6i)$}
\sq{Simplify if $z_{1}=2+3i$ and $z_{2}=4+7i$, then find $z_{1}z_{2}$.}
\sq{If $z=-7-4i$, then show that $z+0=z$.}
\sq{Verify the multiplicative identity for $z=3-2i$}
\sq{Find the additive inverse of $z=7-10i$}
\sq{Find multiplicative inverse of $4-3i$}
\sq{Find the modulus of the following complex number: $4+3i$}
\sq{Find the modulus of the following complex number: $-5-4i$}
\sq{Find the modulus of the following complex number: $\dfrac{3}{5}-\dfrac{4}{5}i$}
\sq{Find the modulus of the following complex number: $-\sqrt{2}-\sqrt{3}\,i$}
\sq{If $z=5-2i$, then verify that: $\overline{z}=z$}
\sq{If $z=5-2i$, then verify that: $|z|=|\overline{z}|$}
\sq{If $z=5-2i$, then verify that: $|z|=|-z|$}

\sq{If $z=5-2i$, then verify that: $z\overline{z}=|z|^{2}$}
\sq{If $z=5-2i$, then verify that: $|z|=|-\overline{z}|$}
\sq{If $z_{1}=2+3i,z_{2}=-1+i$, then evaluate: $Re(z_{1}z_{2})$}
\sq{If $z_{1}=2+3i,z_{2}=-1+i$, then evaluate: $Im(z_{1}z_{2})$}
\sq{For any two complex number $z_{1}$ and $z_{2}$ we have $\overline{z_{1}+z_{2}}=\overline{z_{1}}+\overline{z_{2}}$.}
\sq{If $z_{1}=4+3i,\ z_{2}=5+2i$, then prove that $\overline{z_{1}+z_{2}}=\overline{z_{1}}+\overline{z_{2}}$}
\sq{If $z_{1}=4+3i,\ z_{2}=5+2i$, then prove that $\overline{z_{1}+z_{2}}=\overline{z_{1}}+\overline{z_{2}}$}
\sq{Solve: $\overline{z}=z$}
\sq{If $z=7+3i$, then prove that $\overline{z}=z$.}
\sq{Solve: $|z|=|\overline{z}|$}
\sq{If $z=5+4i$, then show that $|z|=|\overline{z}|$.}
\sq{Solve: $z\overline{z}=|z|^{2}$}
\sq{If $z=5+3i$, then show that $z\overline{z}=|z|^{2}$.}
\sq{Find the real and imaginary parts of the following complex number:\\ $(8-3i)^{2}$}
\sq{Find the real and imaginary parts of the following complex number:\\ $(5+3i)^{-1}$}
\sq{Find the real and imaginary parts of the following complex number:\\ $(4-5i)^{-1}$}
\sq{Find the real and imaginary parts of the following complex number:\\ $(4-3i)^{-2}$}
\sq{Find the real and imaginary parts of the following complex number:\\ $\left(\dfrac{3+2i}{4+3i}\right)^{-1}$}
\sq{Is 0 a complex number? Explain.}
\sq{What is the result of mutiplying a complex number by its conjugate?}
\sq{State the condition for two complex numbers to be equal.}
\sq{Simplify: $i^{37}$}
\sq{Simplify: $i^{13}\times i^{11}$}
\sq{Simplify: $(-i)^{-9}$}
\sq{Simplify: $(3-4i)(5-6i)$}
\sq{Simplify: $(3+4i)\div(5-7i)$}
\sq{Find additive and multiplicative inverse of $z=8+9i$}
\sq{If $z_{1}=5+4i,\ z_{2}=3+2i$, then find: $z_{1}z_{2}$}
\sq{If $z_{1}=5+4i,\ z_{2}=3+2i$, then find: $\dfrac{z_{1}}{z_{2}}$}
\sq{If $z_{1}=5+4i,\ z_{2}=3+2i$, then find: $\overline{z_{1}}\,\overline{z_{2}}$}
\sq{If $z_{1}=5+4i,\ z_{2}=3+2i$, then find: $|z_{1}z_{2}|$}
\sq{Find real and imaginary parts of $z=(2+7i)^{-1}$}
\sq{Write the following quadratic equation in standard form:\\ $3x-1=2x^{2}$}
\sq{Write the following quadratic equation in standard form:\\ $2x(x+1)=4(2x+3)$}
\sq{Write the following quadratic equation in standard form:\\ $2x^{2}-4x=4x+7$}
\sq{Write the following quadratic equation in standard form:\\ $4(3x-2)=9x^{2}$}
\sq{Write the following quadratic equation in standard form:\\ $2x+\dfrac{1}{x}=5-\dfrac{1}{x},\ x\neq0$}
\sq{Write the following quadratic equation in standard form:\\ $\dfrac{6x+6}{20-x}=\dfrac{1}{x},\ x\neq0,20$}
\sq{Solve the following quadratic equation by factorization method:\\ $x^{2}-x-6=0$}
\sq{Solve the following quadratic equation by factorization method:\\ $x^{2}+3x-28=0$}
\sq{Solve the following quadratic equation by factorization method:\\ $6x^{2}+13x-5=0$}
\sq{Use quadratic formula to solve the following equation:\\ $2x^{2}-5x+3=0$}
\sq{Use quadratic formula to solve the following equation:\\ $2x^{2}-7x-15=0$}
\sq{Use quadratic formula to solve the following equation:\\ $2x^{2}+7x=15$}
\sq{Use quadratic formula to solve the following equation: $x^{2}+11=7x$}
\sq{Define algebraic expression.}
\sq{Define linear equation.}
\sq{Define pure quadratic equation.}
\sq{Define fractional equation.}
\sq{Define parabola.}
\sq{Define transposition.}
\sq{Solve by factorization method: $x^{2}+\dfrac{7}{3}x=2$}
\sq{Solve by using quadratic formula: $3x^{2}+7x-6=0$}
\sq{Solve the quadratic equation $y=2x^{2}-x-10=0$}
\sq{Find the points of intersection of the following linear equation with coordinate axes graphically: $x+y=8$}
\sq{Find the points of intersection of the following linear equation with coordinate axes graphically: $x-y=1$}
\sq{Find the points of intersection of the following linear equation with coordinate axes graphically: $x-2y=1$}
\sq{Find the points of intersection of the following linear equation with coordinate axes graphically: $x-2y+2=0$}

\sq{Find the points of intersection of the following linear equation with coordinate axes graphically: $5x-5y=1$}
\sq{Find the points of intersection of $y=2x+4$ with coordinate axes graphically.}
\sq{From a quadratic equation whose roots are given below: -4,9}
\sq{From a quadratic equation whose roots are gven below: 5,-7}
\sq{From a quadratic equation whose roots are given below: $\dfrac{-7}{5}$, $\dfrac{-6}{5}$}
\sq{From a quadratic equation whose roots are given below: $\dfrac{-3}{4}$, $\dfrac{7}{2}$}
\sq{From a quadratic equation whose roots are given below:\\ $3+\sqrt{5}$, $3-\sqrt{5}$}
\sq{From a quadratic equation whose roots are given below:\\ $-2+\sqrt{3}$, $-2-\sqrt{3}$}
\sq{If $\alpha,\beta$ are the roots of the equation $x^{2}+2x+4=0$, then find the equation whose roots are: $\dfrac{1}{\alpha},\dfrac{1}{\beta}$}
\sq{If $\alpha,\beta$ are the roots of the equation $x^{2}+2x+4=0$, then find the equation whose roots are: $\dfrac{\alpha}{\beta},\dfrac{\beta}{\alpha}$}
\sq{Find the sum and the product of the roots of the quadratic equation\\ $x^{2}+6x+8=0$}
\sq{Form the equation whose roots are $\dfrac{2}{3}$ and $\dfrac{3}{4}$}
\sq{Find the condition that the roots of $ax^{2}+bx+c=0$ may be equal in magnitude but opposite in sign, where a is not equal to 0.}
\sq{Examine the nature of roots of the following quadratic equation:\\ $3x^{2}-9x-2=0$}
\sq{Examine the nature of roots of the following quadratic equation:\\ $x^{2}+6x+9=0$}
\sq{Examine the nature of roots of the following quadratic equation:\\ $2x^{2}+4x+5=0$}
\sq{Examine the nature of roots of the following quadratic equation:\\ $7x^{2}-6x-1=0$}
\sq{Examine the nature of roots of the following quadratic equation:\\ $5x^{2}-2x+10=0$}
\sq{Examine the nature of roots of the following quadratic equation:\\ $x^{2}-8x+16=0$}
\sq{Examine the nature of roots of the following quadratic equation:\\ $3x^{2}-8x-2=0$}
\sq{Examine the nature of roots of the following quadratic equation:\\ $4x^{2}+4x+1=0$}
\sq{Examine the nature of roots of the following quadratic equation:\\ $2x^{2}+3x+4=0$}
\sq{Examine the nature of roots of the following quadratic equation:\\ $x^{2}+4x+3=0$}
\sq{Make F the subject of the formula, $C^{o}=\dfrac{5}{9}(F^{o}-32)$}
\sq{The formula for finding simple interest is $I=PRT$: Make P the subject of the formula.}
\sq{The formula for finding simple interest is $I=PRT$: Make T the subject of the formula.}
\sq{Make a the subject of the formula $S=2a+(n-1)d$}
\sq{The volume of a cylinder is given by the formula, $V=\pi r^{2}h$. Make h the subject of the formula.}
\sq{The area of trapezoid is $A=\dfrac{1}{2}h(b_{1}+b_{2})$, make h the subject of the formula.}
\sq{If $y=mx+c$, then make x the subject of this equation.}
\sq{Perimeter (P) of a rectangle is P=2(L+w), make L as the subject of this formula.}
\sq{The equation of a parabola is $y^{2}=4ax$, make x as a subject of this equation.}
\sq{If $P=S-C$ where S is selling price and C is cost price. Make S as subject of the equation.}
\sq{Volume of the cone is $V=\dfrac{1}{3}\pi r^{2}h$, make h as subject of this formula.}
\sq{The area of a circle is $A=\pi r^{2}$. Make r the subject of the formula.}
\sq{A company models its profit P in thousand of rupees by the equation: $P(x)=-5x^{2}+150x-1000$ where x is the price per item in rupees. Find the price that gives maximum profit.}
\sq{A ball is thrown upward with an initial velocity of $40ms^{-1}$. Calculate the maximum height it reaches above ground level.}
\sq{Form a quadratic equation whose roots are: $6,\dfrac{3}{2}$}
\sq{Examine the nature of the roots of the following equation:\\ $15x^{2}+11x+2=0$}
\sq{Examine the nature of roots of the following equation:\\ $x^{2}-x-1=0$}
\sq{Write the number of rows and number of columns in the following matrix. $A=\begin{bmatrix}0&0\\0&0\end{bmatrix}$}
\sq{Write the number of rows and number of columns in the following matrix. $B=\begin{bmatrix}3\\4\\7\end{bmatrix}$}

\sq{Write the number of rows and number of columns in the following matrix. $C=\begin{bmatrix}8&-10&11\end{bmatrix}$}
\sq{Write the number of rows and number of columns in the following matrix. $D=\begin{bmatrix}3&2&1\\1&0&3\\4&5&6\end{bmatrix}$}
\sq{Write the number of rows and number of columns in the following matrix. $E=\begin{bmatrix}5&9&-2\\-3&4&5\end{bmatrix}$}
\sq{Write the order of the matrix: $A=\begin{bmatrix}1\\0\end{bmatrix}$}
\sq{Write the order of the matrix: $B=\begin{bmatrix}3&4\end{bmatrix}$}
\sq{Write the order of the matrix: $C=\begin{bmatrix}3&5\\5&6\end{bmatrix}$}
\sq{Write the order of the matrix: $D=\begin{bmatrix}0&3&4\\5&0&9\end{bmatrix}$}
\sq{Write the order of the matrix: $E=\begin{bmatrix}0&1\\3&4\\5&6\end{bmatrix}$}
\sq{Write the order of the matrix: $F=\begin{bmatrix}5\end{bmatrix}$}
\sq{Is the following matrix is equal? $A=\begin{bmatrix}0&3\\4&1\end{bmatrix}$}
\sq{Is the following matrix is equal? $B=\begin{bmatrix}9\\-7\end{bmatrix}$}
\sq{Is the following matrix is equal? $C=\begin{bmatrix}2\times3&2-1\\2\times2&4-2\\4+4&3+0\end{bmatrix}$}
\sq{Is the following matrix is equal? $D=\begin{bmatrix}5+4\\-8+1\end{bmatrix}$}
\sq{Is the following matrix is equal? $E=\begin{bmatrix}6&1\\4&2\\8&3\end{bmatrix}$}
\sq{Is the following matrix is equal? $F=\begin{bmatrix}3-3&3\\3+1&1\end{bmatrix}$}
\sq{If $\begin{bmatrix}2x+1&4\\7&5\end{bmatrix}=\begin{bmatrix}9&4\\7&y\end{bmatrix}$, then find the values of x and y.}
\sq{Define matrices.}
\sq{Define row matrix.}
\sq{Define column matrix.}
\sq{Define order of matrix.}
\sq{Define equal matrix.}
\sq{Is the following matric is unit matric, row matrices, column matrices or null matrices. $A=\begin{bmatrix}5&7&8\end{bmatrix}$}
\sq{Is the following matric is unit matric, row matrices, column matrices or null matrices. $B=\begin{bmatrix}0\end{bmatrix}$}
\sq{Is the following matric is unit matric, row matrices, column matrices or null matrices. $C=\begin{bmatrix}1&0\\0&1\end{bmatrix}$}
\sq{Is the following matric is unit matric, row matrices, column matrices or null matrices. $D=\begin{bmatrix}0&0\\0&0\end{bmatrix}$}
\sq{Is the following matric is unit matric, row matrices, column matrices or null matrices. $E=\begin{bmatrix}6\\0\\8\end{bmatrix}$}
\sq{Is the following matric is unit matric, row matrices, column matrices or null matrices. $F=\begin{bmatrix}7\\1\\9\end{bmatrix}$}
\sq{Identify type of given matrices as row, coulumn, square and rectangular matrices: $A=\begin{bmatrix}5\\3\\4\end{bmatrix}$}
\sq{Identify type of given matrices as row, coulumn, square and rectangular matrices: $B=\begin{bmatrix}1&6\\4&1\end{bmatrix}$}
\sq{Identify type of given matrices as row, coulumn, square and rectangular matrices: $C=\begin{bmatrix}3&5&8\\0&4&-2\end{bmatrix}$}
\sq{Identify type of given matrices as row, coulumn, square and rectangular matrices: $D=\begin{bmatrix}5&-5\\2&7\end{bmatrix}$}
\sq{Identify type of given matrices as row, coulumn, square and rectangular matrices: $E=\begin{bmatrix}3&2\\4&1\\5&0\end{bmatrix}$}
\sq{Identify type of given matrices as row, coulumn, square and rectangular matrices: $F=\begin{bmatrix}5&-3&7\end{bmatrix}$}
\sq{Identify type of given matrices as row, coulumn, square and rectangular matrices: $G=\begin{bmatrix}3&0&1\\1&3&4\\5&2&-3\end{bmatrix}$}

\sq{Identify type of given matrices as row, coulumn, square and rectangular matrices: $H=\begin{bmatrix}3&5\\4&4\\5&2\end{bmatrix}$}
\sq{Is the following matric is diagonal, scalar or unit matric. $A=\begin{bmatrix}1&0\\0&1\end{bmatrix}$}
\sq{Is the following matric is diagonal, scalar or unit matric. $B=\begin{bmatrix}5&0\\0&5\end{bmatrix}$}
\sq{Is the following matric is diagonal, scalar or unit matric. $C=\begin{bmatrix}3&0\\0&0\end{bmatrix}$}
\sq{Is the following matric is diagonal, scalar or unit matric. $D=\begin{bmatrix}7&0\\0&6\end{bmatrix}$}
\sq{Is the following matric is diagonal, scalar or unit matric. $E=\begin{bmatrix}2&0\\0&2\end{bmatrix}$}
\sq{Find transpose of the following matric: $A=\begin{bmatrix}2&3\\4&5\end{bmatrix}$}
\sq{Find transpose of the following matric: $B=\begin{bmatrix}3\\7\\6\end{bmatrix}$}
\sq{Find transpose of the following matric: $C=\begin{bmatrix}5&-2&4\end{bmatrix}$}
\sq{Find transpose of the following matric: $D=\begin{bmatrix}2&5\\3&6\\4&7\end{bmatrix}$}
\sq{Find negative of the following matric: $A=\begin{bmatrix}-3&0\\5&6\end{bmatrix}$}
\sq{Find negative of the following matric: $B=\begin{bmatrix}-3&3\\-2&2\end{bmatrix}$}
\sq{Find negative of the following matric: $C=\begin{bmatrix}-9&1\\1&-7\end{bmatrix}$}
\sq{If $A=\begin{bmatrix}2&3\\4&5\end{bmatrix}$ and $B=\begin{bmatrix}7&6\\5&8\end{bmatrix}$, then verify that: $(A^{t})^{t}=A$}
\sq{If $A=\begin{bmatrix}2&3\\4&5\end{bmatrix}$ and $B=\begin{bmatrix}7&6\\5&8\end{bmatrix}$, then verify that: $(B^{t})^{t}=B$}
\sq{Define zero or null matrix.}
\sq{Define square matrix.}
\sq{Define diagonal matrix.}
\sq{Define rectangular matrix.}
\sq{Define scalar matrix.}
\sq{Define identity matrix.}
\sq{Define transpose of a matrix.}
\sq{Define symmetric matrix.}
\sq{Define skew-symmetric matrix.}
\sq{Define negative of a matrix.}
\sq{If the following matrix is conformable for addition and subtraction? $A=\begin{bmatrix}3&4\\-2&1\end{bmatrix}$}
\sq{If the following matrix is conformable for addition and subtraction? $B=\begin{bmatrix}3\\4\end{bmatrix}$}
\sq{If the following matrix is conformable for addition and subtraction? $C=\begin{bmatrix}5&2\end{bmatrix}$}
\sq{If the following matrix is conformable for addition and subtraction? $D=\begin{bmatrix}1&7\\2&5\end{bmatrix}$}
\sq{If the following matrix is conformable for addition and subtraction? $E=\begin{bmatrix}2\end{bmatrix}$}
\sq{If the following matrix is conformable for addition and subtraction? $F=\begin{bmatrix}7&11\end{bmatrix}$}
\sq{If the following matrix is conformable for addition and subtraction? $G=\begin{bmatrix}a\\b\end{bmatrix}$}
\sq{If the following matrix is conformable for addition and subtraction? $H=\begin{bmatrix}3\end{bmatrix}$}
\sq{If the following matrix is conformable for addition and subtraction? $M=\begin{bmatrix}l\\m\end{bmatrix}$}
\sq{If $X=\begin{bmatrix}1&-1\\-2&2\end{bmatrix},Y=\begin{bmatrix}1&2\\3&4\end{bmatrix}$ and $Z=\begin{bmatrix}3&0\\0&-2\end{bmatrix}$ then find the following: $X+Y$}
\sq{If $X=\begin{bmatrix}1&-1\\-2&2\end{bmatrix},Y=\begin{bmatrix}1&2\\3&4\end{bmatrix}$ and $Z=\begin{bmatrix}3&0\\0&-2\end{bmatrix}$ then find the following: $Y+7Z$}
\sq{If $X=\begin{bmatrix}1&-1\\-2&2\end{bmatrix},Y=\begin{bmatrix}1&2\\3&4\end{bmatrix}$ and $Z=\begin{bmatrix}3&0\\0&-2\end{bmatrix}$ then find the following: $4X-Z$}
\sq{If $X=\begin{bmatrix}1&-1\\-2&2\end{bmatrix},Y=\begin{bmatrix}1&2\\3&4\end{bmatrix}$ and $Z=\begin{bmatrix}3&0\\0&-2\end{bmatrix}$ then find the following: $X+2y+3Z$}
\sq{If $X=\begin{bmatrix}1&-1\\-2&2\end{bmatrix},Y=\begin{bmatrix}1&2\\3&4\end{bmatrix}$ and $Z=\begin{bmatrix}3&0\\0&-2\end{bmatrix}$ then find the following: $X-4y+Z$}
\sq{If $X=\begin{bmatrix}1&-1\\-2&2\end{bmatrix},Y=\begin{bmatrix}1&2\\3&4\end{bmatrix}$ and $Z=\begin{bmatrix}3&0\\0&-2\end{bmatrix}$ then find the following: $Z-Z$}
\sq{Find the additive inverse of the following matrix: $P=\begin{bmatrix}5\\-7\end{bmatrix}$}
\sq{Find the additive inverse of the following matrix: $Q=\begin{bmatrix}9&-3\end{bmatrix}$}

\sq{Find the additive inverse of the following matrix: $R=\begin{bmatrix}2&-1\\-2&3\end{bmatrix}$}
\sq{Find the additive inverse of the following matrix: $S=\begin{bmatrix}1&0\\0&1\end{bmatrix}$}
\sq{If $A=\begin{bmatrix}2&4\\1&5\end{bmatrix},B=\begin{bmatrix}3&4&1\\5&2&7\end{bmatrix}$ then find $A+B$ and $B-A$, if possible.}
\sq{If $\begin{bmatrix}x&1\\y&2\end{bmatrix}\begin{bmatrix}1&0\\3&-1\end{bmatrix}=\begin{bmatrix}7&-1\\4&-2\end{bmatrix}$ then find the values of x and y.}
\sq{Find AB and BA, if possible: $A=\begin{bmatrix}1&2\\-1&0\end{bmatrix},B=\begin{bmatrix}3&2\\1&-1\end{bmatrix}$}
\sq{Find AB and BA, if possible: $A=\begin{bmatrix}1&-2\end{bmatrix},B=\begin{bmatrix}3\\-4\end{bmatrix}$}
\sq{Find AB and BA, if possible: $A=\begin{bmatrix}4\\4\end{bmatrix},B=\begin{bmatrix}2&5\end{bmatrix}$}
\sq{Find AB and BA, if possible: $A=\begin{bmatrix}1&2\\-1&1\\3&0\end{bmatrix},B=\begin{bmatrix}-1&4&1\\2&3&1\end{bmatrix}$}
\sq{Find AB and BA, if possible: $A=\begin{bmatrix}1&2&2\\3&1&1\end{bmatrix},B=\begin{bmatrix}1&5\\2&4\\1&6\end{bmatrix}$}
\sq{If $\begin{bmatrix}4&a\\b&3\end{bmatrix}\begin{bmatrix}6\\6\end{bmatrix}=\begin{bmatrix}6\\3\end{bmatrix}$ then find the values of and b.}
\sq{If $A=\begin{bmatrix}a&b\\c&d\end{bmatrix},B=\begin{bmatrix}l&m\\n&p\end{bmatrix}$, then find: AB, if possible}
\sq{If $A=\begin{bmatrix}a&b\\c&d\end{bmatrix},B=\begin{bmatrix}l&m\\n&p\end{bmatrix}$, then find: BA if possible.}
\sq{If $A=\begin{bmatrix}2&0\\5&-3\end{bmatrix},B=\begin{bmatrix}-2\\1\end{bmatrix}$ then find: AB, if possible.}
\sq{If $A=\begin{bmatrix}2&0\\5&-3\end{bmatrix},B=\begin{bmatrix}-2\\1\end{bmatrix}$ then find: BA, if possible}
\sq{Find the value of the following determinant: $\begin{vmatrix}10&5\\4&6\end{vmatrix}$}
\sq{Find the value of the following determinant: $\begin{vmatrix}-5&8\\-3&-7\end{vmatrix}$}
\sq{Find the value of the following determinant: $\begin{vmatrix}3&8\\0&2\end{vmatrix}$}
\sq{Find whether the following matrix are singular or non-singular: $A=\begin{bmatrix}5&3\\3&2\end{bmatrix}$}
\sq{Find whether the following matrix are singular or non-singular: $B=\begin{bmatrix}7&21\\2&6\end{bmatrix}$}
\sq{Find whether the following matrix are singular or non-singular: $C=\begin{bmatrix}13&5\\7&3\end{bmatrix}$}
\sq{Find whether the following matrix are singular or non-singular: $D=\begin{bmatrix}2&0\\0&0\end{bmatrix}$}
\sq{Find the value of x when $A=\begin{bmatrix}x&6\\5&15\end{bmatrix}$ is a singular matrix.}
\sq{Find the adjoint of the following matrix: $A=\begin{bmatrix}a&b\\c&d\end{bmatrix},B\begin{bmatrix}5&-2\\3&7\end{bmatrix},C\begin{bmatrix}-3&5\\-3&2\end{bmatrix}$}
\sq{Find the adjoint of the following matrix: $A=\begin{bmatrix}a&b\\c&d\end{bmatrix},B\begin{bmatrix}5&-2\\3&7\end{bmatrix},C\begin{bmatrix}-3&5\\-3&2\end{bmatrix}$}
\sq{Find the adjoint of the following matrix: $A=\begin{bmatrix}a&b\\c&d\end{bmatrix},B\begin{bmatrix}5&-2\\3&7\end{bmatrix},C\begin{bmatrix}-3&5\\-3&2\end{bmatrix}$}
\sq{Find multiplicative inverse of the following matrix: $\begin{bmatrix}5&0\\0&5\end{bmatrix}$}
\sq{Find multiplicative inverse of the following matrix: $\begin{bmatrix}-4&8\\7&2\end{bmatrix}$}
\sq{Find multiplicative inverse of the following matrix: $\begin{bmatrix}40&8\\5&2\end{bmatrix}$}
\sq{Find multiplicative inverse of the following matrix: $\begin{bmatrix}3&5\\5&-3\end{bmatrix}$}
\sq{Find multiplicative inverse of the following matrix: $\begin{bmatrix}10&8\\3&3\end{bmatrix}$}
\sq{Find multiplicative inverse of the following matrix: $\begin{bmatrix}-2&-3\\4&5\end{bmatrix}$}
\sq{Define singular matrix.}
\sq{Define non-singular martix.}
\sq{Define adjoint of matrix.}
\sq{Define multiplicative invers.}
\sq{Find the value of determinant of $B=\begin{bmatrix}8&2\\4&3\end{bmatrix}$.}
\sq{If $A=\begin{bmatrix}x&8\\5&10\end{bmatrix}$ is a singular matrix, then find the value of x.}
\sq{If $A\begin{bmatrix}4&2\\7&6\end{bmatrix},B=\begin{bmatrix}5&5\\8&8\end{bmatrix}$, then find: $(A-B)^{t}$}
\sq{If $A\begin{bmatrix}4&2\\7&6\end{bmatrix},B=\begin{bmatrix}5&5\\8&8\end{bmatrix}$, then find: $B^{t}-A^{t}$}
\sq{If $A\begin{bmatrix}4&2\\7&6\end{bmatrix},B=\begin{bmatrix}5&5\\8&8\end{bmatrix}$, then find: $2A+3B$}
\sq{If $A=\begin{bmatrix}7&3\\2&1\end{bmatrix},B=\begin{bmatrix}5&3\\2&4\end{bmatrix}$, then find: $|B|$}
\sq{If $A=\begin{bmatrix}7&3\\2&1\end{bmatrix},B=\begin{bmatrix}5&3\\2&4\end{bmatrix}$, then find: Adj B}

\sq{If $A=\begin{bmatrix}7&3\\2&1\end{bmatrix},B=\begin{bmatrix}5&3\\2&4\end{bmatrix}$, then find: $A^{-1}$}
\sq{If $A=\begin{bmatrix}7&3\\2&1\end{bmatrix},B=\begin{bmatrix}5&3\\2&4\end{bmatrix}$, then find: $A^{-1}A$}
\sq{If $A=\begin{bmatrix}7&3\\2&1\end{bmatrix},B=\begin{bmatrix}5&3\\2&4\end{bmatrix}$, then find: $(AB)^{t}$}
\sq{If $A=\begin{bmatrix}7&3\\2&1\end{bmatrix},B=\begin{bmatrix}5&3\\2&4\end{bmatrix}$, then find: $(B^{t})^{t}$}
\sq{If $f(x)=2x+5,g(x)=3x-2$, then find: $f(x)+g(x)$}
\sq{If $f(x)=2x+5,g(x)=3x-2$, then find: $f(x)-g(x)$}
\sq{If $f(x)=x+2,\ g(x)=2x+4$, then find: $f(x).g(x)$}
\sq{If $f(x)=x+2,\ g(x)=2x+4$, then find: $g(x).f(x)$}
\sq{If $f(x)=x+2,\ g(x)=2x+4$, then find: $\dfrac{f(x)}{g(x)}$}
\sq{If $f(x)=x+2,\ g(x)=2x+4$, then find: $\dfrac{g(x)}{f(x)}$}
\sq{For the functions f and g, find: $(fog)(x)$ where\\ $f(x)=2x+3\ ;\ g(x)=x^{3}$}
\sq{For the functions f and g, find: $(fog)(x)$ where\\ $f(x)=\dfrac{2}{x},x\neq0\ ;\ g(x)=2x^{2}-1$}
\sq{For the functions f and g, find: $(fog)(x)$ where\\ $f(x)=2x-1\ ;\ g(x)=\dfrac{x+1}{2}$}
\sq{For the functions f and g, find: $(gof)(x)$ where\\ $f(x)=2x+3\ ;\ g(x)=x^{3}$}
\sq{For the functions f and g, find: $(gof)(x)$ where\\ $f(x)=\dfrac{2}{x},x\neq0\ ;\ g(x)=2x^{2}-1$}
\sq{For the functions f and g, find: $(gof)(x)$ where\\ $f(x)=2x-1\ ;\ g(x)=\dfrac{x+1}{2}$}
\sq{For the functions f and g, find: $(fof)(x)$ where\\ $f(x)=2x+3\ ;\ g(x)=x^{3}$}
\sq{For the functions f and g, find: $(fof)(x)$ where\\ $f(x)=\dfrac{2}{x},x\neq0\ ;\ g(x)=2x^{2}-1$}
\sq{For the functions f and g, find: $(fof)(x)$ where\\ $f(x)=2x-1\ ;\ g(x)=\dfrac{x+1}{2}$}
\sq{For the functions f and g, find: $(gog)(x)$ where\\ $f(x)=2x+3\ ;\ g(x)=x^{3}$}
\sq{For the functions f and g, find: $(gog)(x)$ where\\ $f(x)=2x+3\ ;\ g(x)=x^{3}$}
\sq{For the functions f and g, find: $(gog)(x)$ where\\ $f(x)=\dfrac{2}{x},x\neq0\ ;\ g(x)=2x^{2}-1$}
\sq{For the functions f and g, find: $(gog)(x)$ where\\ $f(x)=2x-1\ ;\ g(x)=\dfrac{x+1}{2}$}
\sq{Given that $f(x)=3x+2$ and $g(x)=2x+3$. Find: $g(f(4))$}
\sq{Given that $f(x)=3x+2$ and $g(x)=2x+3$. Find: $f(f(3))$}
\sq{Given that $f(x)=3x+2$ and $g(x)=2x+3$. Find: $f(g(-2))$}
\sq{Find $f^{-1}(x)$ in the following: $f(x)=2x-3$}
\sq{Find $f^{-1}(x)$ in the following: $f(x)=4x^{3}-1$}
\sq{Find $f^{-1}(x)$ in the following: $f(x)=\sqrt{x-1},\ x\geq1$}
\sq{Find $f^{-1}(x)$ in the following: $f(x)=\dfrac{x+1}{3x-2},x\neq\dfrac{2}{3}$}
\sq{The functions f and g are defined such that $f(x)=4x+2$ and $g(x)=6x-18$. find $f^{-1}(x)$ and $g^{-1}(x)$.}
\sq{The functions f and g are defined such that $f(x)=4x+2$ and $g(x)=6x-18$. find x if $f^{-1}(x)$ and $g^{-1}(x)$.}
\sq{Without finding $f^{-1}(x)$, find the domain and range of $f^{-1}(x)$:\\ $f(x)=12x-3$}
\sq{Without finding $f^{-1}(x)$, find the domain and range of $f^{-1}(x)$:\\ $f(x)=\dfrac{1}{2}x+8$}
\sq{Without finding $f^{-1}(x)$, find the domain and range of $f^{-1}(x)$:\\ $f(x)=\dfrac{x}{1+x},x\neq-1$}
\sq{Without finding $f^{-1}(x)$, find the domain and range of $f^{-1}(x)$:\\ $f(x)=\sqrt{x-2},x\geq2$}
\sq{Define composition of the function.}
\sq{Define one-to-one function.}
\sq{Define onto function.}
\sq{Define bijective function.}
\sq{Define inverse function.}
\sq{If $f(x)=x^{2}+3x+1$, then evaluate: $f(0)$}
\sq{If $f(x)=x^{2}+3x+1$, then evaluate: $f(-2)$}
\sq{If $f(x)=x^{2}+3x+1$, then evaluate: $f(a)$}
\sq{If $f(x)=x^{2}+3x+1$, then evaluate: $f(x+1)$.}
\sq{Find the domain and range of $f(x)=(x+1)^{2}$}
\sq{Let $f(x)=2x+3$ and $g(x)=x^{2}+x+1$. Add $f(x)$ and $g(x)$}
\sq{Let $f(x)=5x+2,\ g(x)=x^{2}-3$. Find $f(x)-g(x)$.}

\sq{Let $f(x)=x+1,\ g(x)=2x-3$, find $f(x).g(x)$.}
\sq{Let $f(x)=x^{2}-4,\ g(x)=x+2$. Find $\left(\dfrac{f}{g}\right)(x)$, where $x\neq-2$.}
\sq{If $f(x)=2x+3,\ g(x)=x^{2}$, then find: $(fog)(x)$.}
\sq{If $f(x)=2x+3,\ g(x)=x^{2}$, then find: $(gof)(x)$.}
\sq{If $f(x)=2x+3,\ g(x)=x^{2}$, then find: $(fof)(x)$.}
\sq{If $f(x)=2x+3,\ g(x)=x^{2}$, then find: $(gog)(x)$.}
\sq{Find the inverse of $f(x)=2x+3$. Also find the domain and range of $f^{-1}(x)$.}
\sq{Solve and express the solution on number line: $|x-2|=6$}
\sq{Solve and express the solution on number line: $|2x+1|=7$}
\sq{Solve and express the solution on number line: $|4x-9|=3$}
\sq{Solve and express the solution on number line: $|7-2x|=1$}
\sq{Solve and express the solution on number line: $|5x+8|\leq3$}
\sq{Solve and express the solution on number line: $|4x-12|\leq0$}
\sq{Solve and express the solution on number line: $|3-4x|>0$}
\sq{Solve and express the solution on number line: $|3x-2|\leq12$}
\sq{Solve and express the solution on number line: $|1-2x|>5$}
\sq{Define absolute valued linear equation.}
\sq{Define absolute linear inequality.}
\sq{Plot the graph of the following function: $f(x)=|x|+2$}
\sq{Plot the graph of the following function: $g(x)=|x-3|$}
\sq{Solve $2|x-3|=16$ and express the solution on number line.}
\sq{Solve the inequality $|x-1|\geq5$ and exress the solution on number line.}
\sq{Express the solution of $40|2x-5|+1>201$ on number line.}
\sq{If $f(x)=25-x^{2}$ and $g(x)=5+x$, then find: $f(x)+g(x)$}
\sq{If $f(x)=25-x^{2}$ and $g(x)=5+x$, then find: $f(x)-g(x)$}
\sq{If $f(x)=25-x^{2}$ and $g(x)=5+x$, then find: $f(x).g(x)$}
\sq{If $f(x)=25-x^{2}$ and $g(x)=5+x$, then find: $\dfrac{f(x)}{g(x)}$}
\sq{If $f(x)=25-x^{2}$ and $g(x)=5+x$, then find: $f(7)$}
\sq{If $f(x)=25-x^{2}$ and $g(x)=5+x$, then find: $g(-8)$}
\sq{If $f(x)=x^{3}$ and $g(x)=14+2x$, then find: $(fog)(x)$}
\sq{If $f(x)=x^{3}$ and $g(x)=14+2x$, then find: $(gof)(x)$}
\sq{If $f(x)=x^{3}$ and $g(x)=14+2x$, then find: $(fof)(x)$}
\sq{If $f(x)=x^{3}$ and $g(x)=14+2x$, then find: $(gog)(x)$}
\sq{Reduce the following rational expression to lowest forms:\\ $\dfrac{3a^{2}-6ab}{2a^{2}b-4ab^{2}}$}
\sq{Reduce the following rational expression to lowest forms:\\ $\dfrac{abx+bx^{2}}{acx+cx^{2}}$}
\sq{Reduce the following rational expression to lowest forms:\\ $\dfrac{ac}{a^{2}x^{2}-ax}$}
\sq{Reduce the following rational expression to lowest forms:\\ $\dfrac{15a^{2}b^{2}c}{100(a^{2}-a^{2}b)}$}
\sq{Reduce the following rational expression to lowest forms:\\ $\dfrac{4x^{2}-9y^{2}}{4x^{2}+6xy}$}
\sq{Reduce the following rational expression to lowest forms:\\ $\dfrac{20(x^{3}-y^{3})}{5x^{2}+5xy+5y^{2}}$}
\sq{Reduce the following rational expression to lowest forms:\\ $\dfrac{x(2a^{2}-3ax)}{a(4a^{2}x-9x^{3})}$}
\sq{Reduce the following rational expression to lowest forms:\\ $\dfrac{x^{2}-5x}{x^{2}-4x-5}$}
\sq{Reduce the following rational expression to lowest forms:\\ $\dfrac{3x^{2}+6x}{x^{2}+4x+4}$}
\sq{Reduce the following rational expression to lowest forms:\\ $\dfrac{x^{2}+xy-2y^{2}}{x^{3}-y^{3}}$}
\sq{Reduce the following rational expression to lowest forms:\\ $\dfrac{2x^{2}+17x+21}{3x^{2}+26x+35}$}
\sq{Define algebraic fraction.}
\sq{Define rational expression.}
\sq{Reduce to lowest form $\dfrac{24a^{3}c^{2}x^{2}}{18a^{3}x^{2}-12a^{2}x^{3}}$.}
\sq{Reduce to lowest form: $\dfrac{6x^{2}-8xy}{9xy-12y^{2}}$}
\sq{Simplify: $\dfrac{3}{x-y}+\dfrac{1}{y-x}$}
\sq{Simplify: $\dfrac{x^{3}}{x-y}+\dfrac{y^{3}}{y-x}$}
\sq{Define manipulation.}
\sq{Simplify: $\dfrac{24lm}{5l+10m}\times\dfrac{5l^{2}-20m^{2}}{16mn}$}

\sq{Simplify: $\dfrac{x^{2}-3x+2}{x^{2}+3x-4}\times\dfrac{2x^{2}+8x}{3x+6}$}
\sq{Simplify: $\dfrac{x^{3}-8}{x^{2}-4}\div\dfrac{x^{2}+2x+4}{x^{2}+4x+4}$}
\sq{Simplify: $\dfrac{(a^{2}+ab)^{2}}{(a^{2}-ab)^{2}}\div\left(\dfrac{a+b}{a-b}\right)^{2}$}
\sq{Simplify: $\dfrac{x^{2}-5x+6}{x^{2}-25}\times\dfrac{x^{2}+5x}{2x-6}$}
\sq{Simplify: $\dfrac{3a^{2}-6ab}{6ab}\times\dfrac{10b^{2}}{4ab-8b^{2}}$}
\sq{Simplify: $\dfrac{2a^{2}-2c^{2}}{15a^{2}c^{2}}\div\dfrac{3a+3c}{5ac}$}
\sq{Define exponential equation.}
\sq{Define reciprocal equation.}
\sq{Reduce the following to the lowest form: $\dfrac{x^{2}-7x+12}{x^{2}-6x+9}$}
\sq{Reduce the following to the lowest form: $\dfrac{x^{2}+3x-18}{7x-21}$}
\sq{Name the quadrant in which each point lies: (4,3)}
\sq{Name the quadrant in which each point lies: (5,-4)}
\sq{Name the quadrant in which each point lies: (-6,2)}
\sq{Name the quadrant in which each point lies: (-4,-4)}
\sq{Plot the following points on the coordinate plane: A(3,-3)}
\sq{Plot the following points on the coordinate plane: B(-3,3)}
\sq{Plot the following points on the coordinate plane: C(5,7)}
\sq{Plot the following points on the coordinate plane: D(-2,-4)}
\sq{Name the tail and tip of the following vector: (image)}
\sq{Name the tail and tip of the following vector: (image)}
\sq{Name the tail and tip of the following vector: (image)}
\sq{Name the tail and tip of the following vector: (image)}
\sq{Name the tail and tip of the following vector: (image)}
\sq{Name the tail and tip of the following vector: (image)}
\sq{Write the vector AB in the form of $x\underline{i}+y\underline{j}$:\\ $A(1,-7),B(-2,4)$}
\sq{Write the vector AB in the form of $x\underline{i}+y\underline{j}$: (equation).\\ $A(8,9),B(12,3)$}
\sq{Find the magnitude of the $\underline{a}$: $\underline{a}=-3\underline{i}+2\underline{j}$}
\sq{Find the magnitude of the $\underline{a}$: $\underline{a}=4\underline{i}-3\underline{j}$}
\sq{Find the magnitude of the $\underline{a}$: $\underline{a}=\dfrac{1}{2}\underline{i}+\dfrac{3}{2}\underline{j}$}
\sq{Find a unit vector in the direction of the vector given below:\\ $\underline{a}=-4\underline{i}+5\underline{j}$}
\sq{Find a unit vector in the direction of the vector given below:\\ $\underline{a}=6\underline{i}+8\underline{j}$}
\sq{Find a unit vector in the direction of the vector given below:\\ $\underline{a}=\dfrac{1}{\sqrt{6}}\underline{i}+\dfrac{1}{\sqrt{6}}\underline{j}$}
\sq{Find a unit vector in the direction of the vector given below:\\ $\underline{a}=\dfrac{1}{2}\underline{i}-\dfrac{3}{4}\underline{j}$}
\sq{Is the following vector is parallel? $\underline{a}=6\underline{i}+\underline{j},\ \underline{b}=12\underline{i}+2\underline{j}$}
\sq{Is the following vector is parallel? $\underline{a}=-2\underline{i}+3\underline{j},\ \underline{b}=6\underline{i}-9\underline{j}$}
\sq{Is the following vector is parallel? $\underline{a}=5\underline{i}-4\underline{j},\ \underline{b}=6\underline{i}-3\underline{j}$}
\sq{Is the following vector is parallel? $\underline{a}=3\underline{i}-7\underline{j},\ \underline{b}=6\underline{i}-14\underline{j}$}
\sq{Find a vector thrice in length of $3\underline{i}-2\underline{j}$ but opposite the direction.}
\sq{Define origin.}
\sq{Define scalar quantity.}
\sq{Define vector quantity.}
\sq{Define negative of vector.}
\sq{Define eqaul vector.}
\sq{Define null vector.}
\sq{Define position vector.}
\sq{Define paraller vector.}
\sq{Define unit vector.}
\sq{Write the coordinates of the following points from the given figure: (image)}
\sq{If $P(2,4)$ and $Q(-5,3)$, then find: (equation)}
\sq{If P(2,4) and Q(-5,3), then find: (equation)}
\sq{If equation and equation, then find the following vectors: $\underline{a}+\underline{b}$}
\sq{If equation and equation, then find the following vector: $\underline{a}+3\underline{b}$}
\sq{If equation and equation, then find the following vector: $3\underline{a}+\dfrac{1}{2}\underline{b}$}
\sq{If equation and equation , then find the following vector: $\underline{b}-\underline{a}$}
\sq{If equation and equation, the find the followng vector: $4\underline{b}-5\underline{a}$}
\sq{If equation and equation, then find the following vector: $\dfrac{3}{2}\underline{a}-\underline{b}$}
\sq{If $\underline{a}=6\underline{i}-\underline{j},\ \underline{b}=\underline{i}+5\underline{j}$ and $\underline{c}=3\underline{i}+5\underline{j}$ then find the magnitude of the following vector: $\underline{b}-\underline{c}$}
\sq{If $\underline{a}=6\underline{i}-\underline{j},\ \underline{b}=\underline{i}+5\underline{j}$ and $\underline{c}=3\underline{i}+5\underline{j}$ then find the magnitude of the following vector: $\underline{a}-2\underline{b}+\underline{c}$}
\sq{If $\underline{a}=6\underline{i}-\underline{j},\ \underline{b}=\underline{i}+5\underline{j}$ and $\underline{c}=3\underline{i}+5\underline{j}$, then find the magnitude of the following vector: $\underline{c}-\underline{b}-\underline{a}$}
\sq{Find the values of x and y in the following equation:\\ $(x\underline{i}+y\underline{j})+(2\underline{i}+3\underline{j})=7\underline{i}+6\underline{j}$}

\sq{Find the value of x and y in the following equation:\\ $(x\underline{i}-5\underline{j})+(3\underline{i}+5\underline{j})=8\underline{i}+y\underline{j}$}
\sq{Find the value of x and y in the following equation:\\ $(y\underline{i}+3\underline{j})+(-5\underline{i}+2x\underline{j})=9\underline{i}+7\underline{j}$}
\sq{If $\underline{a}=-2\underline{i}+7\underline{j}$ and $\underline{b}=3\underline{i}-5\underline{j}$ then find the vector components of $\underline{a}+5\underline{b}$.}
\sq{What is parallelogram law of addition.}
\sq{Define coplanar vectors.}
\sq{If $\underline{a}=2\underline{i}+3\underline{j},\ \underline{b}=4\underline{i}-6\underline{j}$ then find $\underline{a}+\underline{b}$}
\sq{If $\underline{a}=6\underline{i}+13\underline{j}$ and $\underline{b}=-4\underline{i}+7\underline{j}$, then find $\underline{a}-\underline{b}$.}
\sq{If $\underline{a}=2\underline{i}-3\underline{j}$ and $\underline{b}=3\underline{i}+4\underline{j}$ then find the vector components of $3\underline{a}+4\underline{b}$.}
\sq{A car enters a loop with velocity vector $\underline{v}=30\underline{j}$ and exits with velocity $\underline{v}'=30\underline{i}$ What is the change in velocity vector?}
\sq{Define translation vector.}
\sq{Define isosceles triangle.}
\sq{Calculate area of the following triangle ABC.\\ $a=7\ cm,\ b=8\ cm,\ \gamma=38^{o}$}
\sq{Calculate area of the following triangle ABC.\\ $a=11\ cm,\ c=14\ cm,\ \beta=51^{o}$}
\sq{Calculate area of the following triangle ABC.\\ $b=3\ cm,\ c=9\ cm,\ \alpha=78^{o}$}
\sq{Calculate area of the following triangle ABC.\\ $a=10\ cm,\ \alpha=62^{o},\ \beta=69^{o}$}
\sq{Calculate area of the following triangle ABC.\\ $c=4\ cm,\ \beta=36^{o},\ \gamma=80^{o}$}
\sq{Calculate area of the following triangle ABC.\\ $c=6.6\ cm,\ \alpha=23^{o},\ \gamma=89^{o}$}
\sq{Calculate area of the following triangle ABC.\\ $a=5.3\ cm,\ b=4.7\ cm,\ c=8.2\ cm$}
\sq{Calculate area of the following triangle ABC.\\ $a=6.12\ cm,\ b=8.34\ cm,\ c=7.12\ cm$}
\sq{Define oblique triangles.}
\sq{What is sine law.}
\sq{What cosine law}
\sq{Calculate area of triangle ABC in which $a=4.2,\ b=3.2$ and $\gamma=70^{o}$.}
\sq{Calculate area of $\triangle ABC$ in which\\ $a=4\ cm,\ b=3\ cm,\ c=5\ cm$}
\sq{Calculate area of triangle ABC in which\\ $\alpha=35^{o},\ \beta=65^{o},\ a=2cm$.}
\sq{A boat crosses a river 80 m wide, at a $60^{o}$ angle to the current. What distance does the boat actually travel?}
\sq{In a sound lab, two microphones are placed c=3 m apart and a sound source is a=2 m from one and b=2.5 m from the other . Find the angle between the sound paths a and b.}
\sq{A plane takes off at an angle of elevation of $12^{o}$ and travels 2000 metres along this path. How high is the plane after traveling this distance?}
\sq{Find bearing of point P in the following:\pfig{f461.png}{4.5cm}}
\sq{Find bearing of point P in the following:\pfig{f462.png}{4.5cm}}
\sq{Find bearing of point P in the following:\pfig{f463.png}{4.5cm}}
\sq{Find bearing of point P in the following:\pfig{f464.png}{4.5cm}}
\sq{Calculate the area of $\triangle ABC$, in which:\\ $a=4\ cm,\ b=6\ cm,\ c=8\ cm$}
\sq{Calculate the area of $\triangle ABC$, in which:\\ $b=2.1\ cm,\ c=5\ cm,\ \gamma=45^{o}$}
\sq{Calculate the area of $\triangle ABC$, in which:\\ $c=3.1\ cm,\ \gamma=44^{o},\ \alpha=36^{o}$}
\sq{A chord 8 cm long is at a distance of 3 cm from the centre. calculate the radius of the circle.}

\sq{In the given figure $m\overline{CD}=121$ units and $m\overline{BC}=3x$ units. Find the value of x.\pfig{f469.png}{4.2cm}}
\sq{Define circle.}
\sq{Define circumfrence.}
\sq{Define radial segment.}
\sq{Define diameter.}
\sq{Define chord.}
\sq{Define semi-circle.}
\sq{Define major arc.}
\sq{Define minor arc.}
\sq{Define segment.}
\sq{Define sector.}
\sq{Define angle.}
\sq{Define central angle.}
\sq{Define secant.}
\sq{Define arc.}
\sq{Prove that a unique circle can pass through the points A(1,1), B(4,2), C(3,5).}
\sq{In a circular fountain of radius 8 m, A pipe AB spans across the fountain. If $m\overline{AB}=12\ m$ and O is centre of the circle. Find the distance from O to the chord AB.}
\sq{Two equa-lengths ropes (chords) are tied in a circular playground. One is 8 m from the centre. Where should the second be placed?}
\sq{In a circle with centre at F, $m\overline{NP}=12\ cm$ and $m\overline{EF}=m\overline{FM}$. Find the length of DE.}
\sq{In a round table design, two decorations are placed symmetrically 4 cm from the centre. One chord is 14 cm. What is the length of the other?}
\sq{In the given figure, if $m\overline{PQ}=m\overline{QR}=m\overline{RS}=m\overline{ST}$ and $m\angle POT=70^{o}$, then find the value of x.\pfig{f489.png}{4.2cm}}
\sq{In the given figure, find the values of x and y.\pfig{f490.png}{4.2cm}}
\sq{Find the values of x and y in the given figure, such that $m\widehat{AB}=m\widehat{BD}$.\pfig{f491.png}{4.5cm}}
\sq{In the given, find the value of x.\pfig{f492.png}{4.2cm}}
\sq{Find the value of x in the given figure.\pfig{f493.png}{4.5cm}}
\sq{In the given figure, $\overline{PT}$ is a diameter of circle. if $m\overline{PQ}=m\overline{QR}=m\overline{RS}=m\overline{ST}$, then find a.\pfig{f494.png}{4.5cm}}

\sq{A chord 6 cm long is at a distance of 4 cm from the centre. Calculate the radius of the circle.}
\sq{Find the value of x\pfig{f496.png}{3.8cm}}
\sq{Find the angle ABC\pfig{f497.png}{3.8cm}}
\sq{D,F and F are point on the circumference of a circle with centre at O and $m\angle DOF=130^{o}$. Find $m\angle DEF$.}
\sq{X,Y and Z are points on the circumference of a circle with centre at O and $m\angle XYZ=85^{o}$. Find x.}
\sq{The radii of the two circles are 4cm and 6cm. What is the distance between their centres of they touch externally?}
\sq{An inner decorative circle touches the outer clock face internally. If the outer circle has radius 12 cm and the inner radius 7.5 cm, find the distance between their centres.}
\sq{Find the angles x and y in the following figure: (image)}
\sq{Find the angles x and y in the following figure: (image)}
\sq{Find the angles x and y in the following figure. (image)}
\sq{Find the angle x in the given figure. (image)}
\sq{A circular garden has a walking path forming a quadrilateral inscribed in it. If one angle is $87^{o}$, what is the opposite angle?}
\sq{Find $\theta$ when: $\ell=3\ cm,\ r=2.2\ cm$}
\sq{Find $\theta$, when: $\ell=5.6\ cm,\ r=2\ cm$}
\sq{Find r when: $\ell=5.5\ cm,\theta=40^{o}20'$}
\sq{Find r when: $\ell=13\ cm,\theta=70^{o}$}
\sq{Find $\ell$ and area of sector, when: $r=1.7\ cm,\ \theta=0.25$}
\sq{Find $\ell$ and area of sector, when: $r=3\ cm,\ \theta=45^{o}$}
\sq{Uzam cut a pizza of radius 14 cm into 8 equal slices. What is the area of one slice (sector).}
\sq{S,T,U and V are points on the circumference of a circle. Find the angles x and y. (image)}
\sq{Find the area of sector of a circle having central angle $60^{o}$ and radius 7 cm. Also find the length of the arc.}
\sq{The length of an arc and area of sector of a circle are 5cm and $25cm^{2}$ respectively. Find radius and the central angle of sector.}
\sq{A triangular frame within a semicircular gate joins the ends of the diameter to the peak. What is the angle formed at the peak?}
\sq{In the adjoining figure, find the radius of the circle. (image)}
\sq{Find the angles a,b,c and d in the given figure, if O is the centre of the circle. (image)}
\sq{Find the angles a,b and c if $\overrightarrow{ML}$ is a tangent line and O is the centre of the circle. (image)}
\sq{Find the angles marked in the given figure, where O is the centre of the circle. (image)}
\sq{Construct a circle with the help of given radius and verify its centre by construction: r = 1.5 cm.}
\sq{Construct a circle with the help of given radius and verify its centre by construction. r=1.7 cm.}
\sq{Construct a circle with the help of given radius and verify its centre by construction: r=2cm.}
\sq{Take any three non-collinear P,Q,R and construct a cricle passing through these points.}
\sq{Draw an arc ABC and complete a circle by finding its centre.}
\sq{Draw an arc PQR and complete a circle without finding its centre.}
\sq{Draw tangent to $\widehat{APB}$ at point P, when P is midpoint of the arc.}
\sq{Draw tangent to $\widehat{PQR}$ at point P, when P is endpoint of the arc.}
\sq{Draw a circle of radius 1.3 cm and draw a tangent at point P, when P lies on its circumference.}
\sq{Draw a circle of radius 1.5cm and draw a tangent at point P, when P is at a distance of 8 cm from its centre.}
\sq{Draw a tangent to any point on a circular part of the track having radius = 2cm.}
\sq{Draw a circle of radius 1.4 cm and draw a tangent at point P, when P lies on its circumference.}
\sq{Draw a circle of radius 1.2cm and draw a tangent at point P, when P is at a distance of 5cm from the centre.}
\sq{Draw a circle with radius 1.5 cm. Draw two tangents to this circle that meet an angle of $30^{o}$.}
\sq{Find $Q_{1},Q_{3}$ median, minimum and maximum values for the following box-and-whisker plot. (image)}

\sq{Draw the box-and-whisker plot, of min. value = 12, max. value = 26, median= 19, $Q_{1}=14$ and $Q_{3}=24$.}
\sq{Construct a scatter diagram and draw a line of best fit for the following quiz scores for 8 students in a class:\par\smallskip\noindent\hspace*{7em}\begin{tabular}{|l|c|c|c|c|c|c|c|c|}\hline \textbf{1st quiz}&14&28&34&36&43&45&46&49\\ \textbf{2nd quiz}&16&20&33&35&38&39&42&43\\ \hline\end{tabular}\par\smallskip\hspace*{7em}Describe correlation between 1st and 2nd quiz scores also.}
\sq{Define cumulative frequency.}
\sq{What is table of cumulative frequency.}
\sq{Define cumulative frequency polygon.}
\sq{Define median.}
\sq{Define quartile.}
\sq{Deinfe decile.}
\sq{Define percentiles.}
\sq{Define interquartile range.}
\sq{What box and whisker plot?}
\sq{What is Box?}
\sq{What is whisker?}
\sq{Define scatter diagram.}
\sq{Define a line of best fit.}
\sq{Define correlation.}
\sq{Define positive correlation.}
\sq{Define negative correlation.}
\sq{Define zero correlation.}
\sq{Defne no-correlation.}
\sq{For the given set of data: 2,10,3,16,14,13,4,7,11,5,17,14,18 What is the median of the data set?}
\sq{For the given set of data: 2,10,3,16,14,13,4,7,11,5,17,14,18 What are the upper and lower quartiles?}
\sq{For the given set of data: 2,10,3,16,14,13,4,7,11,5,17,14,18. What are the extreme values of the data set?}
\sq{For the given set of data:2,10,3,16,14,13,4,7,11,5,17,14,18 Draw a box-and-whisker plot to display the data.}
\sq{Find the range of the following data set: 63,89,98,125,79,108,117,60}
\sq{Find the range of the following data set: 43.5,13.6,18.9,38.4,61.4,29.4}
\sq{If the range and the lowest value of a set of data are 46.7 and 13.4 respectively, then find the highest value.}
\sq{Calculate the range of the following:\par\smallskip\noindent\hspace*{7em}\begin{tabular}{|l|c|c|c|c|c|}\hline \textbf{Income (in Rs.)}&4000-4500&4500-5000&5000-5500&5500-6000&6000\\ \textbf{No. of workers}&8&12&30&21&6\\ \hline\end{tabular}}
\sq{A company forecasts monthly sales (rupees in millions): 15,18,14,20,13. Find the variability in sales predictions.}
\sq{Define measures of dispersion.}
\sq{Define range.}
\sq{Define variance.}
\sq{Define standar deviation.}
\sq{Define mean deviation.}
\sq{Define quartile deviation.}
\sq{Define coefficient of variations.}
\sq{Find the range of the following data: 15,22,18,30,17,25}
\sq{The number of carpets sold in each day of a week is 13,8,4,9,7,12,10. Find the variance.}
\sq{The summary statistics for a data set is given below. Show it with a box-and-whisker plot.\par\smallskip\noindent\hspace*{7em}\begin{tabular}{|c|c|c|c|c|}\hline Min. Value&Max. Value&Q1&Median&Q3\\ 5&70&12.6&43&55.6\\ \hline\end{tabular}}
\sq{Find the range of the given date: 25,30,35,40,50,60,65,75.}
\sq{Find the sample space for tossing three coins using three diagram.}
\sq{If A is an event of a random experiment such that $P(A):P(\overline{A})=17:15$ and n (S)=640, then find (i) $P(\overline{A})$ (ii) $n(A)$}
\sq{A coin is tossed thrice. What is the probability of getting two consecutive tails?}
\sq{If the probability of an event A is $\dfrac{7}{10}$, then find the probability of the event Not A.}
\sq{What is ranom experiment.}
\sq{Define outcomes.}
\sq{Define sample space.}
\sq{Define event.}
\sq{Define equally likely events.}
\sq{Define certain events.}
\sq{Define impossible events.}
\sq{Define mutually exclusive events.}
\sq{Define exhaustive events.}
\sq{Define complementry events.}

\sq{Define elementry events.}
\sq{Define combined events.}
\sq{Write the name of types of combined events.}
\sq{Define idependent events}
\sq{Define dependent events.}
\sq{Define non-mutually exclusive events.}
\sq{Define tree diagram.}
\sq{Define sample space diagram.}
\sq{Define possibility diagram.}
\sq{Define probability of an event.Write its formula.}
\sq{Express the sample space for rolling two dice using tree diagram.}
\sq{Hashim rolled a dice and tossed a coin simultaneously. Draw possibility diagram.}
\sq{Two coind are tossed together. What is the probability of getting different faces on the coins?}
\sq{If $P(A)=0.37,\ P(B)=0.42,\ P(A\cap B)=0.09$ then find $P(A\cup B)$.}

\end{document}
